\problem{1}
На плоскости фиксирован остроугольный треугольник $A B C$ с наибольшей стороной $B C$. Пусть $P Q$ - произвольный диаметр его описанной окружности, причём точка $P$ лежит на меньшей дуге $A B$, а точка $Q$ - на меньшей дуге $A C$. Точки $X, Y$ и $Z$ - основания перпендикуляров, опущенных из точки $P$ на прямую $A B$, из точки $Q$ на прямую $A C$ и из точки $A$ на прямую $P Q$. Докажите, что центр описанной окружности треугольника $X Y Z$ лежит на фиксированной окружности (не зависящей от выбора точек $P$ и $Q$ ).

\textbf{Решение.}
Пусть описанная окружность треугольника $A B C$ является единичной с центром в 0 , и треугольник $A B C$ положительно ориентирован. Заметим, что тогда $p=-q$. Обозначим угол $\angle B A C$ через $\varphi$. Введем обозначение $\zeta=\cos \varphi+i \sin \varphi$. Тогда $\angle P O C=2 \varphi$, откуда $c / b=\zeta^2$. Вычисляем $x=\frac{a+b+p-\bar{p} a b}{2}, y=\frac{a+c-p+\bar{p} a c}{2}$.
Обозначим через $K(k)$ центр описанной окружности треугольника $X Y Z$. Заметим, что $\angle P A B+\angle C A Q=90^{\circ}-\varphi$. Тогда

$$
\angle X Z Y=180^{\circ}-\angle P Z X-\angle Y Z Q=180^{\circ}-\angle P A B-\angle C A Q=90^{\circ}+\varphi
$$


Получаем, что ориентированный угол $X K Y$ равен $180^{\circ}+2 \varphi$. То есть $(y-k)=-\zeta^2(x-k)=(k-x) c / b$, откуда находим

$$
k=\frac{b y+c x}{b+c}=\frac{b a+b c-b p+\bar{p} a b c+c a+c b+c p-\bar{p} a b c}{2(b+c)}=\frac{(c-b)}{2(b+c)} p+\frac{a b+a c+2 b c}{2(b+c)}
$$


То есть координата центра описанной окружности треугольника $X Y Z$ имеет вид $u p+v$, где $u, v$ не зависят от $p$, а $P$ бегает по единичной окружности. Тогда понятно, что и $K$ бегает по окружности, полученной из единичной умножением на $u$ и сдвигом на $v$.


\problem{2}
В треугольнике $A B C$ биссектрисы $A A_1$ и $C C_1$ пересекаются в точке $I$. Прямая, проходящая через точку $B$ параллельно $A C$ пересекает лучи $A A_1$ и $C C_1$ в точках $A_2$ и $C_2$ соответственно. Точка $O_a$ - центр описанной окружности треугольника $A C_1 C_2$, точка $O_c$ - центр описанной окружности треугольника $C A_1 A_2$. Докажите, что $\angle O_a B O_c=\angle A I C$.

\textbf{Решение.}
Пусть описанная окружность треугольника $A B C$ является единичной с центром в нуле. Обозначим через $k$ комплексное число, отвечающее повороту на $\angle B A C=\alpha$ против часовой стрелки, через $l$ - отвечающее повороту на $\angle B C A=\gamma$ по часовой стрелке. Тогда центр вписанной окружности имеет координату $t=b k+b l-b k l$. Обозначим середины дуг $B C$ и $A B$ через $A^{\prime}$ и $C^{\prime}$ соответственно. Тогда $a^{\prime}=b k, c^{\prime}=b l$. Найдем координату точки $c_2$. Во-первых $\left.\left(b-c_2\right) \overline{(a-c)}\right)=\overline{\left(b-c_2\right)}(a-c)$, откуда $c_2+\overline{c_2} b^2 k^2 l^2=b+b k^2 l^2$. При этом $C_2$ лежит на хорде $C C^{\prime}$, откуда $c_2+\overline{c_2} b^2 k^2 l=b k^2+b l$. Решая полученную систему, находим $c_2=\frac{b k^2 l+b l^2-b-b k^2 l^2}{l-1}=b l+b-b k^2 l$. Аналогично $a_2=b k+b-b l^2 k$. Заметим, что $\angle A C_1 C_2=2 \alpha+\gamma$. Тогда ориентированный угол $A O_a C_2=4 \alpha+2 \gamma$, откуда $\left(a-o_a\right) \cdot k^2 / l=\left(c_2-o_a\right)$. Итого, $O_a$ имеет комплексную координату

$$
\begin{gathered}
o_a=\frac{a k^2-c_2 l}{k^2-l}=\frac{2 b k^2 l^2-b l^2-b l}{k^2-l} \\
o_a-b=\frac{2 b k^2 l^2-b l^2-b l-b k^2+b l}{k^2-l}=\frac{2 b k^2 l^2-b k^2-b l^2}{k^2-l}
\end{gathered}
$$


Аналогично $o_c-b=\frac{2 b k^2 l^2-b k^2-b l^2}{l^2-k}$. Тогда

$$
\frac{\left(o_a-b\right)(c-t)}{\left(o_c-b\right)(a-t)}=\frac{\left(l^2-k\right)(b k+b)(k+l)}{\left(k^2-l\right)(b l+b)(k+l)}=\frac{\left(l^2-k\right)(k+1)}{\left(k^2-l\right)(l+1)}
$$


Последнее выражение очевидно вещественное, что и требовалось доказать.


\problem{3}
В остроугольном неравнобедренном треугольнике $A B C$ проведены медиана $A M$ и высота $A H$. На прямых $A B$ и $A C$ отмечены точки $Q$ и $P$ соответственно так, что $Q M \perp A C$ и $P M \perp A B$. Описанная окружность треугольника $P M Q$ пересекает прямую $B C$ вторично в точке $X$. Докажите, что $B H=C X$.

\textbf{Решение.}
второе решение. Пусть описанная окружность треугольника $A B C$ является единичной с центром в нуле. Поскольку $P M \perp A B$, имеем $(p-m)=(\bar{p}-\bar{m}) \cdot a b$, откуда $p-\bar{p} a b=m-\bar{m} a b$. С другой стороны $p+\bar{p} a c=a+c$. Решая систему на $p$ и $\bar{p}$, находим $p=\frac{m c-\bar{m} a b c+a b+b c}{b+c}=\frac{3 b c+a b+c^2-a c}{2(b+c)}$. Аналогичное выражение получается для $q=\frac{3 b c+a c+b^2-a b}{2(b+c)}$.
Отметим на прямой $B C$ точку $X$ так, что $\overline{B X}=\overline{H C}$, тогда $x-b=c-h$, откуда $x=b+c-h$. Вспомнив, что $h=\frac{a+b+c-\bar{a} b c}{2}$, находим $x=\frac{b+c-a+\bar{a} b c}{2}$.
Нам достаточно показать, что $X, P, Q, M$ лежат на одной окружности. Для этого нужно посчитать двойное отношение. Сделаем необходимые вычисления

$$
\begin{gathered}
p-x=\frac{3 b c+a b+c^2-a c-b^2-c^2-2 b c+a b-\bar{a} b^2 c-+a c-\bar{a} b c^2}{2 a(b+c)}= \\
=\frac{a b c+2 a^2 b-a b^2-b^2 c-b c^2}{2 a(b+c)}=\frac{b(a+c)(2 a-b-c)}{2 a(b+c)} \\
p-m=\frac{3 b c+a b+c^2-a c-b^2-c^2-2 b c}{2(b+c)}=\frac{b c+a b-a c-b^2}{2(b+c)}=\frac{(a-b)(b-c)}{2(b+c)}
\end{gathered}
$$


Аналогично $q-x=\frac{c(a+b)(2 a-b-c)}{2 a(b+c)}, q-m=\frac{(a-c)(c-b)}{2(b+c)}$. Наконец, можно посчитать двойное отношение

$$
\frac{(p-x)(q-m)}{(q-x)(p-m)}=\frac{b(a+c)(2 a-b-c)(a-c)(c-b)}{c(a+b)(2 a-b-c)(a-b)(b-c)}=-\frac{b(a+c)(a-c)}{c(a+b)(a-b)}
$$


Последнее выражение действительно вещественное, что сразу следует из подстановки $\bar{a} \rightarrow \frac{1}{a}, \bar{b} \rightarrow \frac{1}{b}, \bar{c} \rightarrow \frac{1}{c}$.

